Generative agents are increasingly deployed in open environments populated by other agents, and recent incidents on agent social networks and in evaluation sandboxes show that small groups of agents can coordinate to the detriment of the rest of the system. We study the problem of detecting such \emph{colluding subnetworks} from the outside: an observer sees, at each time step, a collection of responses from each agent to a collection of queries, where the queries may differ across agents, may be sparse, and may themselves be unobserved and require estimation. We formalize the agents as random functions from a query space to a response space, define a colluding subnetwork as a subset of agents whose conditional behavior shifts in a shared direction after a burn-in period, and pose detection and membership estimation as a scan test over a family of candidate subsets. Building on the Data Kernel Perspective Space and its product-kernel generalization, which represents black-box models from sparse and unpaired responses, we embed each (agent, time) pair in a common low-dimensional space, measure each agent's displacement, and scan for a subset with an anomalously large shared displacement. The null distribution is obtained by permuting each agent's query--response pairs across time, which is exact under the null of no behavioral change. We characterize Type~I error, power, and membership recovery in simulation as a function of effect size, the number of agents, the number of responses per agent per time step, the size of the colluding subnetwork, and whether queries are observed or estimated, and we apply the method to observed behavior on the agent social network Moltbook.
